\chapter{Conclusion and Outlook}
\label{ch:conclusion}

\providecommand{\WL}{\textsc{wl}}

This book has built a single arc. It began with the observation that the
architectures behind the deep-learning revolution --- fully connected,
convolutional, recurrent, and attention-based networks --- were all conceived
for data laid out on regular Euclidean domains, and that this assumption breaks
down for the molecules, networks, and surfaces that pervade the sciences. From
there it developed the mathematical language of symmetry
(\Cref{ch:foundations}), turned that language into a constructive design
principle (\Cref{ch:priors}), and followed the principle through the five
recurring themes --- grids, groups, graphs, geodesics, and gauges
(\Cref{ch:grids,ch:groups,ch:graphs,ch:geodesics,ch:gauges}) --- before
surveying what the resulting models achieve in practice (\Cref{ch:applications}).

This closing chapter does not introduce new machinery. Its purpose is to tie the
threads together: to show that the chapters describe one idea seen from several
angles, to be honest about where that idea reaches its present limits, and to
point to the directions in which it is likely to grow.

\section{A unified view}
\label{sec:concl:unified}

The deepest lesson of geometric deep learning is that architectures which look
unrelated are special cases of a single construction. A fully connected layer, a
convolution, a graph network, and a self-attention block differ not in kind but
in three geometric choices: the \emph{domain} on which the signal lives, the
\emph{symmetry group} the architecture is required to respect, and the
\emph{connectivity} through which information is allowed to flow.

The thread that makes this concrete is \emph{message passing}. A message-passing
layer operates on a graph by computing, for every node, a message from each of
its neighbours, aggregating those messages, and updating the node's state. The
remarkable fact, established across the preceding chapters, is that the principal
deep-learning architectures are this same layer with different graphs and
different weight-sharing rules.

\begin{itemize}
  \item A \emph{fully connected layer} is message passing on the complete graph
  with a distinct weight for every pair of features. It assumes no geometric
  structure whatsoever, which is precisely why it has no inductive bias and must
  learn every invariance from data.
  \item A \emph{convolutional layer} (\Cref{ch:grids}) is message passing on a
  regular grid in which the weight depends only on the \emph{relative} position
  of sender and receiver. This weight sharing \emph{is} translation equivariance,
  and it is what reduces the parameter count from quadratic in the signal size to
  linear in the kernel support.
  \item A \emph{self-attention layer} (the transformer) is message passing on the
  complete graph in which the weights are computed dynamically from the content
  of the nodes. The communication pattern is learned per input rather than fixed,
  buying global, adaptive connectivity at quadratic cost in the sequence length.
  \item A \emph{graph neural network} (\Cref{ch:graphs}) is the general case:
  message passing on an arbitrary graph with permutation symmetry, of which the
  three architectures above are the regular, complete, and adaptive specializations.
\end{itemize}

Group-equivariant networks (\Cref{ch:groups}) refine the convolutional case by
replacing translations with a larger global symmetry, so that weight sharing is
governed by the action of a group rather than a grid. Learning on manifolds
(\Cref{ch:geodesics}) replaces the ambient grid with intrinsic distances and
curvature, recovering message passing as the discretization of a continuous
diffusion. Gauge-equivariant networks (\Cref{ch:gauges}) acknowledge that on a
curved domain there is no canonical local frame, and demand that the construction
be independent of that arbitrary choice. In each step the same skeleton ---
local messages, symmetry-constrained weights, equivariant composition --- is
carried to a more general geometry.

This unification is not merely a tidy taxonomy. It carries genuine explanatory
and predictive force. It explains \emph{why} a given architecture suits a given
kind of data: the symmetry of the domain dictates the equivariance constraints
the network must satisfy, and those constraints in turn fix the form of the
admissible layers. And it is predictive: given a new domain, the framework tells
us which symmetries must be respected and, through equivariant universality
results, what expressive guarantees the resulting model can offer. Geometric deep
learning is therefore not only a description of the architectures we already have
but a recipe for designing the ones we do not.

\section{Open problems}
\label{sec:concl:open}

The same theory that unifies these architectures also exposes, with unusual
clarity, the boundaries of what is currently possible. Four limitations stand out
because they emerge from the formalism itself rather than from engineering
contingency.

\paragraph{Expressivity and the Weisfeiler--Leman barrier.}
The expressive power of standard message-passing networks is bounded exactly by
the one-dimensional Weisfeiler--Leman graph-isomorphism test: two graphs that the
test cannot tell apart will receive identical representations from any such
network. Higher-order ($k$-\WL{}) variants break through this ceiling, but their
cost grows exponentially in $k$, which makes them impractical even for graphs of
moderate size. Designing architectures that surpass the one-dimensional barrier
at polynomial cost --- for instance through structural features or higher-order
attention --- remains a central open problem.

\paragraph{Rigidity versus adaptability.}
There is a fundamental tension between baking symmetry into the architecture and
leaving it to be learned. Equivariant convolutional networks encode symmetry
rigidly: they are data-efficient and come with theoretical guarantees, but they
struggle when the symmetry is only approximate. Attention mechanisms are flexible
and learn context-dependent relations, but they demand more data and offer weaker
structural regularization. There is reason to believe the two paradigms converge
under suitable conditions, yet the optimal point on the rigidity--adaptability
spectrum for any concrete domain is not understood.

\paragraph{Gauge equivariance and global topology.}
The gauge formalism (\Cref{ch:gauges}) gives a rigorous account of local
symmetries, but the global topology of the domain resists it. Non-trivial bundles
and non-trivial holonomy around loops pose difficulties that are at once
theoretical and computational, and current gauge-equivariant networks rely on
discretizations that may fail to preserve the global topological properties of
the underlying space. Closing the gap between the elegance of the formalism and
efficient, topology-faithful implementations is a significant practical
challenge.

\paragraph{Higher-order structures.}
The tools that capture relations beyond pairwise interactions --- Hodge
Laplacians, simplicial and cellular complexes, persistent homology --- are
mathematically mature, but the architectures built on them are not. Designing
convolution and attention operations on general cellular complexes that respect
the algebraic structure of the underlying topology without incurring prohibitive
cost is an open and increasingly active problem.

\section{Future directions}
\label{sec:concl:future}

The limitations above suggest, almost by negation, where the field is heading.
Three directions follow directly from the material of this book.

\paragraph{Hybrid convolution--attention architectures.}
The formal correspondence between convolution and attention, together with the
analysis of the rigidity--adaptability spectrum, points to architectures that
combine the two in a controlled way. The goal is to interpolate between the
data efficiency of equivariant convolutions and the expressiveness of attention,
rather than choosing one wholesale. The transformer's success on grids ---
domains where convolution already excels --- shows that the boundary between the
two is porous, and a principled way of placing a model anywhere along it would be
of broad value.

\paragraph{Complex geometric domains.}
The five-Gs framework extends naturally to product spaces that mix curvatures
--- hyperbolic, spherical, and Euclidean --- within a single model, which is
appropriate for data that is simultaneously hierarchical and cyclic. In the
topological direction, the theory of cellular sheaves provides a foundation for
\emph{sheaf neural networks}, which generalize graph networks by equipping the
relations between nodes with the algebraic structure of a sheaf. Both extensions
enlarge the class of domains the framework can address while keeping its
equivariant backbone intact.

\paragraph{Generalized harmonic analysis.}
Harmonic analysis has been a unifying tool throughout, from spectral convolutions
to spherical harmonics. Its central result, the Peter--Weyl theorem, requires the
symmetry group to be compact. Extending the framework to non-compact groups ---
the Lorentz group, the affine groups --- calls for non-commutative harmonic
analysis beyond what the present theory provides, and would widen the reach of
equivariant networks to symmetries that matter in relativistic physics and
general signal processing.

\section{A closing reflection}
\label{sec:concl:reflection}

The argument of this book is that geometry is not a collection of tricks for
squeezing extra accuracy out of a model, but an organizing principle that
explains why particular architectures work for particular data. The differences
between fully connected networks, convolutional networks, graph networks, and
transformers reflect three choices --- domain, symmetry group, and connectivity
--- and once those choices are made explicit, the architectures cease to be a
zoo and become a family.

The mathematics that sustains this view --- Lie theory, Riemannian geometry,
gauge theory, algebraic topology --- has a reach that goes beyond describing the
models we already have. Equivariant universality results and the correspondences
between paradigms give the framework predictive power: they let us anticipate
which geometric constraints a domain demands and what guarantees the resulting
architecture will offer. In that sense geometric deep learning is less a finished
theory than a programme, and the problems left open in this chapter are an
invitation to continue it. As the data we wish to learn from grows in geometric
richness, the case for letting that geometry shape the model --- rather than the
other way around --- only grows stronger~\citep{bronstein2021geometric}.
